\documentclass[11pt]{article}

\usepackage[margin=1in]{geometry}
\usepackage{amsmath,amssymb,amsthm,mathtools}
\usepackage{enumitem}
\usepackage{booktabs}
\usepackage[round,authoryear]{natbib}
\usepackage{hyperref}
\usepackage{algorithm}
\usepackage{algorithmic}

\theoremstyle{plain}
\newtheorem{theorem}{Theorem}
\newtheorem{lemma}[theorem]{Lemma}
\newtheorem{proposition}[theorem]{Proposition}
\newtheorem{corollary}[theorem]{Corollary}

\theoremstyle{definition}
\newtheorem{definition}[theorem]{Definition}
\newtheorem{hypothesis}[theorem]{Hypothesis}

\theoremstyle{remark}
\newtheorem{remark}[theorem]{Remark}

\DeclareMathOperator{\im}{Im}
\DeclareMathOperator{\re}{Re}

\newcommand{\R}{\mathbb{R}}
\newcommand{\C}{\mathbb{C}}
\newcommand{\Q}{\mathbb{Q}}
\newcommand{\SO}{\mathrm{SO}}
\newcommand{\OG}{\mathrm{O}}
\newcommand{\SH}{Y}
\newcommand{\CG}[6]{C^{#5,\,#6}_{#1,\,#2;\,#3,\,#4}}
\newcommand{\abs}[1]{\lvert #1 \rvert}
\newcommand{\norm}[1]{\lVert #1 \rVert}
\newcommand{\conj}[1]{\overline{#1}}
\newcommand{\Sb}{\mathcal{S}_\beta}
\newcommand{\Sseed}{\mathcal{S}_{\mathrm{seed}}}
\newcommand{\SP}{\mathcal{S}_P}
\newcommand{\Phiaug}{\Phi_{\mathrm{aug}}}
\newcommand{\Useed}{\mathcal{U}_{\mathrm{seed}}}

\title{Completeness of a Selective $\SO(3)$-Bispectrum on $S^2$
       \\ (with optional Clebsch--Gordan power diagnostics)}
\author{}
\date{}

\begin{document}
\maketitle

\begin{abstract}
We give a conditional proof that an $\Theta(L^2)$-sized purely scalar
selective invariant $\Phiaug$, consisting of a constant-size full
low-degree bispectrum block plus $O(\ell)$ selected scalar bispectral
entries at each higher degree, separates $\SO(3)$-orbits of generic
real-valued band-limited signals. The two load-bearing assumptions are
isolated explicitly: finite-band injectivity of the full low-degree
scalar seed, and nonvanishing of the per-degree bootstrap
determinants. The proof follows Kakarala's pattern, motivated by
\citet{kakarala1992triple,kakarala2012bispectrum}: a full-bispectrum
seed followed by a triangular recurrence in which each
$\mathbf{F}_\ell$ is recovered by solving an explicit linear system.
Clebsch--Gordan power entries are optional implementation diagnostics
and are not used in the proof; the adjective ``augmented'' refers to
this implementation variant.
\end{abstract}

\section{Setup and notation}
\label{sec:setup}

\paragraph{Signal space.}
Let $V_L$ denote the space of real-valued band-limited functions
$f : S^2 \to \R$ with band-limit~$L$, expanded in spherical harmonics
$\{\SH_\ell^m\}$:
\begin{equation}\label{eq:sh-expansion}
  f(\theta,\varphi)
  = \sum_{\ell=0}^{L}\sum_{m=-\ell}^{\ell}
    a_\ell^m \, \SH_\ell^m(\theta,\varphi),
  \qquad a_\ell^{-m} = (-1)^m \conj{a_\ell^m}
  \quad (\text{reality}).
\end{equation}
Write $\mathbf{F}_\ell := (a_\ell^{-\ell},\dots,a_\ell^{\ell}) \in \C^{2\ell+1}$ for the
coefficient vector at degree~$\ell$. The reality constraint reduces
$\mathbf{F}_\ell$ to $2\ell+1$ real degrees of freedom, so
$\dim_\R V_L = \sum_{\ell=0}^{L}(2\ell+1) = (L+1)^2$.

\paragraph{Group action.}
Under $g\in\SO(3)$, the coefficient vectors transform irreducibly,
$\mathbf{F}_\ell \mapsto D^\ell(g)\,\mathbf{F}_\ell$, where $D^\ell$ is
the $(2\ell+1)$-dimensional Wigner $D$-matrix.

\begin{definition}[Bispectrum on $S^2$]\label{def:bispectrum}
For each triple $(\ell_1,\ell_2,\ell)$ with $0 \le \ell_1,\ell_2,\ell \le L$
satisfying the triangle inequality $\abs{\ell_1{-}\ell_2} \le \ell \le \ell_1{+}\ell_2$,
the \emph{scalar bispectral coefficient} is
\begin{equation}\label{eq:bispectrum}
  \beta_{\ell_1,\ell_2,\ell}(f)
  := \sum_{\substack{m_1,m_2,m \\ m_1+m_2 = m}}
     \CG{\ell_1}{m_1}{\ell_2}{m_2}{\ell}{m} \,
     a_{\ell_1}^{m_1}\, a_{\ell_2}^{m_2}\, \conj{a_\ell^m},
\end{equation}
where $\CG{\ell_1}{m_1}{\ell_2}{m_2}{\ell}{m} := \langle \ell_1, m_1;\, \ell_2, m_2 \mid \ell, m\rangle$
is the Clebsch--Gordan coefficient. Each $\beta_{\ell_1,\ell_2,\ell}$
is $\SO(3)$-invariant. The full bispectrum at band-limit~$L$ has
$\Theta(L^3)$ entries~\citep{kakarala1992triple,kakarala2012bispectrum}.
\end{definition}

\begin{definition}[Clebsch--Gordan power spectrum]\label{def:cg-power}
For a triple $(\ell_1,\ell_2,\ell)$ satisfying the triangle inequality,
let $(\mathbf{F}_{\ell_1}\otimes\mathbf{F}_{\ell_2})\mid_\ell$ denote the
projection of the Kronecker product onto the $D^\ell$-isotypic component,
realised by contracting against the CG coefficients. Define the
\emph{CG power spectrum entry}
\begin{equation}\label{eq:cg-power}
  P_{\ell_1,\ell_2,\ell}(f)
  := \norm{(\mathbf{F}_{\ell_1}\otimes\mathbf{F}_{\ell_2})\mid_\ell}^2
   = \sum_{m=-\ell}^{\ell}
     \Bigl| \sum_{m_1+m_2=m}
       \CG{\ell_1}{m_1}{\ell_2}{m_2}{\ell}{m}\,
       a_{\ell_1}^{m_1}\, a_{\ell_2}^{m_2} \Bigr|^2.
\end{equation}
This is a degree-$4$, real-valued, $\SO(3)$-invariant polynomial in
the signal coefficients.
\end{definition}

\noindent
The standard power spectrum $\norm{\mathbf{F}_\ell}^2$ is the special
case $P_{0,\ell,\ell}/\abs{a_0^0}^2$ (with our CG normalisation). The
CG power generalises it by projecting
$\mathbf{F}_{\ell_1}\otimes\mathbf{F}_{\ell_2}$ into each output
channel~$V_\ell$ before taking norms.

\begin{definition}[Reflection map $T_R$]\label{def:TR}
The azimuthal reflection $R:(\theta,\varphi)\mapsto(\theta,-\varphi)$ acts
on coefficients by $a_\ell^m\mapsto(-1)^m a_\ell^{-m}$. For real signals
this equals complex conjugation: $a_\ell^m\mapsto\conj{a_\ell^m}$. Since
$R\in\OG(3)\setminus\SO(3)$, the signal $T_R\cdot f$ is generically in a
\emph{different} $\SO(3)$-orbit from $f$, and any complete
$\SO(3)$-invariant must distinguish $f$ from $T_R\cdot f$.
\end{definition}

\begin{definition}[Selective invariant]\label{def:Phi-aug}
Let $L_{\mathrm{seed}}:=\min(L,7)$. Define the \emph{seed index set}
\[
  \Sseed := \{(\ell_1,\ell_2,\ell):
    0\le\ell_1\le\ell_2\le L_{\mathrm{seed}},\;
    \abs{\ell_1-\ell_2}\le\ell\le\min(\ell_1+\ell_2,L_{\mathrm{seed}})\},
\]
the \emph{high-degree bootstrap blocks}
$\mathcal{T}_\ell$ ($\ell\ge8$) of~\eqref{eq:Tell}, and the scalar
index set
\[
  \Sb := \Sseed \;\cup\; \bigcup_{\ell=8}^L \mathcal{T}_\ell .
\]
The \emph{selective invariant} $\Phiaug:V_L\to\R^N$ is the
concatenation of all selected scalar bispectral entries
$\{\beta_{\ell_1,\ell_2,\ell}\}_{(\ell_1,\ell_2,\ell)\in\Sb}$, stored
as complex numbers (equivalently, as real and imaginary parts). The
seed sub-block has size at most $(L_{\mathrm{seed}}{+}1)^3\le 512$,
independent of the final band-limit once $L\ge7$. We set
$\SP:=\varnothing$ in the formal theorem.
\end{definition}

\begin{remark}[On the name ``augmented'']\label{rem:augmented-name}
The implementation variant of $\Phiaug$ used in our experiments
appends a CG-power block (Remark~\ref{rem:cg-optional}) and is what we
call \emph{augmented}. The conditional theorem below is proved for the
purely scalar selective invariant ($\SP=\varnothing$); the CG-power
block is optional throughout the proof.
\end{remark}

\begin{remark}[Optional CG-power diagnostics]
\label{rem:cg-optional}
Any $O(L^2)$ family $\{P_{\ell_1,\ell_2,\ell}\}_{\SP}$ of CG-power
entries may be appended to $\Phiaug$ without affecting completeness or
its asymptotic size. These entries are useful implementation
diagnostics (Section~\ref{sec:cgpower}); they are not used by the
conditional theorem below.
\end{remark}

\begin{definition}[Genericity]\label{def:gen}
A signal $f \in V_L$ is \emph{generic} if it satisfies:
\begin{enumerate}[nosep,label=(G\arabic*)]
  \item $a_0^0 \ne 0$ (nonzero mean);
  \item $\norm{\mathbf{F}_1}\ne 0$ (nonzero dipole);
  \item $a_2^1 \ne 0$ after the gauge of Lemma~\ref{lem:gauge}
    (azimuthal phase resolvable);
  \item the gauge-fixed seed truncation belongs to the seed-injective
    stratum $\Useed$ of Hypothesis~\ref{hyp:seed};
  \item $\Delta_\ell \ne 0$ for all $8 \le \ell \le L$, where
    $\Delta_\ell$ is the real-coordinate determinant of
    Hypothesis~\ref{hyp:bootstrap-det}.
\end{enumerate}
Conditions (G1)--(G5) are polynomial inequalities in the gauge-fixed
real coordinates of the gauge slice $\mathcal{G}$ of
Lemma~\ref{lem:gauge} (forward reference) and therefore cut out a
Zariski-open subset $\mathcal{V}\subset\mathcal{G}$. Its
$\SO(3)$-saturation $\SO(3)\cdot\mathcal{V}\subset V_L$ is a generic
co-null subset of $V_L$, whose complement is contained in a finite
union of proper algebraic subvarieties of $\mathcal{G}$ and is in
particular Lebesgue-null.
\end{definition}

\section{The inversion algorithm}
\label{sec:algorithm}

The completeness theorem (Theorem~\ref{thm:main}) is the statement
that the procedure below recovers $f$ from $\Phiaug(f)$ up to the
$\SO(3)$-action for generic~$f$, provided the bootstrap determinants
are nonzero at the required degrees.

\begin{algorithm}[h]
  \caption{Bispectrum inversion on $S^2$ from the selective invariant
    $\Phiaug$ ($\SP=\varnothing$).
    Inputs from $\Phiaug(f)$ are typeset \textbf{in bold}; recovered
    coefficients accumulate in $(\mathbf{F}_0,\dotsc,\mathbf{F}_L)$
    on the gauge slice.}
  \label{alg:so3-inversion}
  \begin{algorithmic}[1]
    \STATE \textbf{Input:} $\Phiaug(f)$, band-limit $L \ge 4$, seed cutoff
      $L_{\mathrm{seed}}=\min(L,7)$.
    \STATE $A_0 \leftarrow (\boldsymbol{\beta_{0,0,0}})^{1/3}$
      \hfill \COMMENT{closed form, Lemma~\ref{lem:F0F1}}
    \STATE $c \leftarrow \sqrt{\boldsymbol{\beta_{0,1,1}}\,/\,A_0}$
      \hfill \COMMENT{$c = \norm{\mathbf{F}_1}$}
    \STATE $\mathbf{F}_0 \leftarrow A_0$,\quad
      $\mathbf{F}_1 \leftarrow (0,\,c,\,0)^\top$
      \hfill \COMMENT{gauge: align $\mathbf{F}_1$ with $\hat z$}
    \STATE $(\mathbf{F}_2,\dots,\mathbf{F}_{L_{\mathrm{seed}}}) \leftarrow
      \textsc{SeedSolve}\bigl(\boldsymbol{\{\beta_{\ell_1,\ell_2,\ell}\}_{\ell_1,\ell_2,\ell\le L_{\mathrm{seed}}}}\bigr)$
      \hfill \COMMENT{returns gauge slice rep with $a_2^1\in\R_{>0}$ (Lemma~\ref{lem:gauge}, Hypothesis~\ref{hyp:seed})}
    \FOR{$\ell = L_{\mathrm{seed}}+1, \dots, L$}
      \STATE Build $A_\ell \in \C^{(2\ell+1)\times(2\ell+1)}$, $b_\ell\in\C^{2\ell+1}$
        from selected bispectral entries linear in $\mathbf{F}_\ell$
        and depending on $(\mathbf{F}_0,\dotsc,\mathbf{F}_{\ell-1})$
      \STATE $\mathbf{F}_\ell \leftarrow A_\ell^{-1} b_\ell$
        \hfill \COMMENT{bootstrap determinant nonzero}
    \ENDFOR
    \STATE \textbf{Return} $(\mathbf{F}_0,\dotsc,\mathbf{F}_L)$, equivalently $f$, up to $\SO(3)$.
  \end{algorithmic}
\end{algorithm}

The algorithm is a direct generalisation of the finite-group selective
inversion of \citet[Algorithm~3]{mataigne2024selective}: the for-loop
in Lines~7--10 plays the role of their breadth-first traversal of the
irrep graph, with the per-step Fourier-coefficient recurrence replaced
by an $O(\ell)$-dimensional linear solve. The seed step in Line~6
plays the role of the trivial-irrep base case in Mataigne et al.\ but
is no longer in closed form; it is the finite-band seed assumption
isolated in Hypothesis~\ref{hyp:seed}.

\paragraph{Index sets.}
For each output degree $\ell\ge8$, the bootstrap block referenced in
Definition~\ref{def:Phi-aug} is
\begin{equation}\label{eq:Tell}
\begin{aligned}
  \mathcal{T}_\ell
  :=&\ \{(a,\ell,\ell-a) : 1\le a\le \ell-1\}\\
   &\cup \{(a,\ell,\ell-a+1) : 2\le a\le \ell-1\}\\
   &\cup \{(a,\ell-a,\ell) : 1\le a\le4\},
\end{aligned}
\end{equation}
with $(\ell-1)+(\ell-2)+4=2\ell+1$ entries. Combined with the seed
block $\Sseed$, this gives $|\Sb|=|\Sseed|+\sum_{\ell=8}^L(2\ell+1)
=O(L^2)$ (Proposition~\ref{prop:size}). Optional CG-power
diagnostics (Remark~\ref{rem:cg-optional}) may be appended separately,
still at $O(L^2)$ cost.

\section{CG-power diagnostics}
\label{sec:cgpower}

The argument of \citet{mataigne2024selective} for finite groups never
needs the CG power: the scalar bispectrum on $L^2(G)$ already has full
real Jacobian rank, because the reality constraint on a finite-group
signal is encoded by a pairing of each irrep $\rho_i$ with its complex
conjugate irrep $\rho_{\bar i}$, which is itself an irrep of $G$.

For $\SO(3)$ on $S^2$, an earlier real-valued implementation of the
selected bispectrum chose only the parity-reduced scalar component of
each entry (real or imaginary part dictated by parity). With that row
selection, the per-degree real Jacobian was rank-deficient at every
degree $\ell\ge4$. CG-power entries close that gap diagnostically and
were originally introduced for that reason.

The proof in this paper does \emph{not} rely on those parity-reduced
real rows. Each entry of $\Sb$ contributes its full complex value to
$\Phiaug$ (Definition~\ref{def:Phi-aug}), and the per-degree bootstrap
of Section~\ref{sec:bootstrap} uses the complex linear system of
Hypothesis~\ref{hyp:bootstrap-det}. By Lemma~\ref{lem:bootstrap},
invertibility of $A_\ell$ over $\C$ already forces uniqueness of
$\mathbf{F}_\ell$ on the real-signal slice, so no real-rank gap arises
in the proof itself. CG-power entries remain useful as implementation
diagnostics and as a richer invariant for downstream learning, but
they are not required for the conditional theorem.

\section{Main theorem}
\label{sec:main}

\begin{theorem}[Conditional completeness]\label{thm:main}
Let $L \ge 4$, and let $f,f'\in V_L$ both be generic in the sense of
Definition~\ref{def:gen}. Assume Hypothesis~\ref{hyp:seed} for
$L_{\mathrm{seed}}=\min(L,7)$ and
Hypothesis~\ref{hyp:bootstrap-det} for every $8\le\ell\le L$. If
$\Phiaug(f)=\Phiaug(f')$, then $f'=g\cdot f$ for some $g\in\SO(3)$.

The output $\Phiaug$ has $\Theta(L^2)$ real components
(Proposition~\ref{prop:size}), matching the orbit-space
dimension lower bound $(L+1)^2 - 3$ up to a constant factor.
\end{theorem}

\noindent
The proof reduces to gauge-fixing (\S\ref{sec:gauge}), the finite
seed hypothesis (\S\ref{sec:seed}), parity bookkeeping
(\S\ref{sec:cgpower-formal}), and the per-degree determinant bootstrap
(\S\ref{sec:bootstrap}). It is collected in \S\ref{sec:proof}.

\subsection{Gauge fixing}
\label{sec:gauge}

\begin{lemma}[Gauge fixing]\label{lem:gauge}
For any $f \in V_L$ satisfying (G1)--(G3) of Definition~\ref{def:gen},
there is a unique $g \in \SO(3)$ such that $\tilde f := g\cdot f$ has
\begin{enumerate}[nosep,label=(\roman*)]
  \item $a_1^{-1} = a_1^1 = 0$ and $a_1^0 = c > 0$
    (alignment of $\mathbf{F}_1$ with $\hat z$), and
  \item $a_2^1 \in \R_{>0}$ (residual $\SO(2)_z$ phase fixed).
\end{enumerate}
This consumes all three rotational degrees of freedom of $\SO(3)$. The
\emph{gauge slice} $\mathcal{G}\subset V_L$ defined by (i)--(ii) has
$\dim_\R\mathcal{G} = (L+1)^2 - 3$.
\end{lemma}
\begin{proof}
$\SO(3)$ acts transitively on $S^2$, so the unit vector
$\mathbf{F}_1/\norm{\mathbf{F}_1}$ can be rotated to $\hat z$; this fixes
$g$ modulo the stabiliser $\SO(2)_z$ of $\hat z$. On the resulting
slice, $\SO(2)_z$ acts on $a_2^1$ by multiplication with $e^{i\varphi}$,
so taking $\varphi = -\arg(a_2^1)$ uniquely places $a_2^1$ on the
positive real axis. The dimension count is
$(L+1)^2 - \dim\SO(3) = (L+1)^2 - 3$.
\end{proof}

\begin{lemma}[$\mathbf{F}_0$ and $\mathbf{F}_1$]\label{lem:F0F1}
On the gauge slice, $\beta_{0,0,0} = (a_0^0)^3$ determines $a_0^0 = A_0$
uniquely (the real cube root), and $\beta_{0,1,1} = A_0\norm{\mathbf{F}_1}^2$
determines $c = \norm{\mathbf{F}_1}$. By (i) of Lemma~\ref{lem:gauge},
$\mathbf{F}_1 = (0,c,0)^\top$.
\end{lemma}

\subsection{Finite-band seed hypothesis}
\label{sec:seed}

\begin{hypothesis}[Finite seed injectivity]\label{hyp:seed}
For each $s\le7$, there is a Zariski-open gauge-fixed seed stratum
$\Useed^{(s)}$ such that the full scalar seed map
\[
  (\mathbf{F}_0,\dots,\mathbf{F}_s)
  \longmapsto
  \{\beta_{\ell_1,\ell_2,\ell} :
    0\le\ell_1\le\ell_2\le s,\;
    \abs{\ell_1-\ell_2}\le\ell\le\min(\ell_1+\ell_2,s)\}
\]
is injective on $\Useed^{(s)}$. In Definition~\ref{def:gen}, $\Useed$
means $\Useed^{(L_{\mathrm{seed}})}$.
\end{hypothesis}

\begin{remark}[Relation to Kakarala]
Kakarala's homogeneous-space completeness theorem
\citep[Theorem~7]{kakarala2012bispectrum}, originally proved in
\citet{kakarala1992triple}, is the template for
Hypothesis~\ref{hyp:seed}: for $H=\SO(2)_z$, it recovers a left
$H$-invariant function on $\SO(3)$ from the full bispectrum under a
maximal $H$-rank condition, up to the normaliser $N_H\subset\SO(3)$.
This does not by itself prove the finite-band seed statement above,
because a band-limited truncation has zero Fourier coefficients above
the cutoff. We therefore use Kakarala only as structural motivation
unless a finite-band version or an independent finite certificate is
provided.
\end{remark}

\begin{remark}[Possible symbolic certificate]
\label{rem:bandeira}
\citet{bandeira2017estimation} give computer-algebra evidence for
generic orbit recovery from polynomial invariants at small band-limits.
To use that result as a proof of Hypothesis~\ref{hyp:seed}, one must
check that its invariant map is exactly the full scalar $S^2$
bispectrum seed above, on the same gauge-fixed real signal space, and
that its certificate establishes generic finite-fibre injectivity
rather than only local rank. We do not assume that identification here.
\end{remark}

\begin{remark}[On the ``up to reflection'' caveat of \citet{edidin2024orbit}]
\label{rem:edidin}
\citet[Section~1]{edidin2024orbit} state that
\citep{kakarala2012bispectrum} recovers the orbit ``up to reflection''.
This caveat is sharp for the constructive recursion in Kakarala's
Section~5 (a polar decomposition step pins $\hat F(1)$ only modulo a
real orthogonal matrix, an $\OG(3)$-ambiguity that propagates), but it
is not present in Theorem~7 itself, which gives recovery up to
$N_H \subset \SO(3)$. The odd-parity entry $\beta_{2,3,4}$ in the
seed block separates $f$ from $T_R\cdot f$
(Corollary~\ref{cor:first-odd}), confirming that no $\OG(3)\setminus\SO(3)$
ambiguity remains.
\end{remark}

\subsection{Parity and the first parity-breaking entry}
\label{sec:cgpower-formal}

\begin{lemma}[Parity of the bispectrum]\label{lem:parity}
For every signal:
\begin{equation}\label{eq:parity-TR}
  \beta_{\ell_1,\ell_2,\ell}(T_R\cdot f)
  = (-1)^{\ell_1+\ell_2+\ell}\, \beta_{\ell_1,\ell_2,\ell}(f).
\end{equation}
For real signals, the further reality constraint
\begin{equation}\label{eq:parity-real}
  \conj{\beta_{\ell_1,\ell_2,\ell}(f)}
  = (-1)^{\ell_1+\ell_2+\ell}\, \beta_{\ell_1,\ell_2,\ell}(f)
\end{equation}
holds. Together: when $\ell_1{+}\ell_2{+}\ell$ is even, $\beta$ is real
and $T_R$-invariant; when odd, $\beta$ is purely imaginary and
$\im\beta$ flips sign under~$T_R$. The CG power $P_{\ell_1,\ell_2,\ell}$
is real and $T_R$-invariant in all cases.
\end{lemma}
\begin{proof}
For~\eqref{eq:parity-TR}, substitute $a_\ell^m \mapsto (-1)^m a_\ell^{-m}$
into~\eqref{eq:bispectrum}; the three sign factors give
$(-1)^{m_1+m_2+m} = (-1)^{2m} = 1$, and re-indexing via the CG symmetry
$\CG{\ell_1}{-m_1}{\ell_2}{-m_2}{\ell}{-m} = (-1)^{\ell_1+\ell_2-\ell}
\CG{\ell_1}{m_1}{\ell_2}{m_2}{\ell}{m}$ yields the claim.
For~\eqref{eq:parity-real}, take complex conjugates inside~\eqref{eq:bispectrum},
use $\conj{a_\ell^m} = (-1)^m a_\ell^{-m}$, and re-index.
For $P$, the CG-projected component picks up
$(-1)^{\ell_1+\ell_2+\ell+M}$ under $T_R$, and re-indexing $M \mapsto -M$
in $\sum_M\abs{\cdot}^2$ leaves the sum unchanged.
\end{proof}

\begin{corollary}[First nonzero odd-parity entry]\label{cor:first-odd}
For real signals, every odd-parity bispectral entry
$\beta_{\ell_1,\ell_2,\ell}$ whose index multiset
$\{\ell_1,\ell_2,\ell\}$ contains a repeated value vanishes
identically; in particular, every odd-parity entry with
$\max(\ell_1,\ell_2,\ell)\le3$ vanishes. The first surviving
odd-parity entry is therefore $\beta_{2,3,4}$
(with $\ell_1+\ell_2+\ell=9$, all distinct). It is generically
nonzero, so $\im\beta_{2,3,4}(f)\ne0$ separates $f$ from $T_R\cdot f$.
\end{corollary}
\begin{proof}
\emph{Repeated-index vanishing.} Define the symmetric scalar tensor
\[
  T_{\ell_1,\ell_2,\ell_3}
  := \sum_{m_1,m_2,m_3}
       \begin{pmatrix} \ell_1 & \ell_2 & \ell_3 \\
                       m_1 & m_2 & m_3 \end{pmatrix}\,
       a_{\ell_1}^{m_1}\, a_{\ell_2}^{m_2}\, a_{\ell_3}^{m_3} .
\]
Converting the Wigner~$3$j into a CG via
$\bigl(\begin{smallmatrix} j_1 & j_2 & j_3 \\
       m_1 & m_2 & m_3\end{smallmatrix}\bigr)
=\frac{(-1)^{j_1-j_2-m_3}}{\sqrt{2j_3+1}}\,
\CG{j_1}{m_1}{j_2}{m_2}{j_3}{-m_3}$ and applying the reality identity
$a_{\ell_3}^{m_3}=(-1)^{m_3}\conj{a_{\ell_3}^{-m_3}}$ shows
$\beta_{\ell_1,\ell_2,\ell}=(-1)^{\ell_1+\ell_2}\sqrt{2\ell+1}\,
T_{\ell_1,\ell_2,\ell}$, so $\beta$ and $T$ have the same vanishing
locus. The 3j symbol is invariant under even column permutations and
picks up $(-1)^{\ell_1+\ell_2+\ell_3}$ under odd ones. If two columns
$i,j$ have $\ell_i=\ell_j$, the inserted vectors $a_{\ell_i},
a_{\ell_j}$ are literally the same coefficient vector of $f$, so
swapping those columns leaves the contracted sum unchanged while
multiplying it by $(-1)^{\ell_1+\ell_2+\ell_3}$. For odd parity this
forces $T=-T$, hence $T=0$ and $\beta=0$.
For $\max(\ell_1,\ell_2,\ell)\le3$, every admissible odd-parity
triangle (e.g.\ $(0,1,2)$, $(1,1,1)$, $(0,3,3)$, $(1,2,3)$, $(2,2,3)$)
has a repeated index, so $\beta_{2,3,4}$ is the first candidate that
survives.

\emph{Nonvanishing of $\beta_{2,3,4}$.} Take the explicit signal
\[
  a_2^2 = a_3^2 = 1, \quad a_4^4 = i,
\]
with all other coefficients in $\mathbf{F}_2,\mathbf{F}_3,\mathbf{F}_4$
set to zero, and the negative-$m$ entries fixed by reality
($a_2^{-2}=a_3^{-2}=1$, $a_4^{-4}=-i$). Only the index choices
$(m_1,m_2,m)\in\{(2,2,4),(-2,-2,-4)\}$ contribute. Using
$\CG{2}{2}{3}{2}{4}{4}=\sqrt{2/5}$ and the standard sign relation
$\CG{2}{-2}{3}{-2}{4}{-4}=-\sqrt{2/5}$,
\[
  \beta_{2,3,4}
  = \sqrt{2/5}\cdot 1\cdot1\cdot(-i)
   + (-\sqrt{2/5})\cdot 1\cdot1\cdot(+i)
  = -2i\sqrt{2/5}\ne0.
\]
Hence $\beta_{2,3,4}$ is not identically zero on $V_L$. Genericity is
then a Zariski-open condition on the gauge slice.
\end{proof}

\subsection{Per-degree linear bootstrap}
\label{sec:bootstrap}

For $\ell\ge8$, use the block $\mathcal{T}_\ell$ in~\eqref{eq:Tell}.
The first two families are cross entries $(a,\ell,q)$, hence linear in
$\mathbf{F}_\ell$:
\[
  \beta_{a,\ell,q}
  = \sum_{m=-\ell}^{\ell} H_{a,q}^{m}\, a_\ell^m,
  \qquad
  H_{a,q}^{m}
  := \sum_{r=-a}^{a}
     \CG{a}{r}{\ell}{m}{q}{r+m}\,
     a_a^r\,\conj{a_q^{r+m}} .
\]
The third family consists of chain entries $(a,\ell-a,\ell)$, linear
in $\conj{\mathbf{F}_\ell}$:
\[
  \beta_{a,\ell-a,\ell}
  = \sum_{m=-\ell}^{\ell} G_a^m\,\conj{a_\ell^m},
  \qquad
  G_a^m
  := \sum_{r+s=m}
     \CG{a}{r}{\ell-a}{s}{\ell}{m}\,
     a_a^r\,a_{\ell-a}^s .
\]
Conjugating the chain equations gives
$\sum_m \conj{G_a^m}a_\ell^m=\conj{\beta_{a,\ell-a,\ell}}$.
Thus all rows in $\mathcal{T}_\ell$ assemble into a square complex
linear system
\begin{equation}\label{eq:bootstrap-system}
  A_\ell(\mathbf{F}_{<\ell})\,\mathbf{F}_\ell=b_\ell .
\end{equation}
Although \eqref{eq:bootstrap-system} is complex-linear in the unknown
$\mathbf{F}_\ell$, its coefficients are functions on the real-signal
lower-degree slice. Whenever a conjugated lower coefficient appears,
we rewrite it using the real coordinates
\[
  a_j^0\in\R,\qquad a_j^m=x_{j,m}+i y_{j,m},\qquad
  a_j^{-m}=(-1)^m(x_{j,m}-i y_{j,m})\quad(m>0).
\]
After this substitution, every entry of $A_\ell$ is a complex-valued
polynomial in the real variables $\{a_j^0,x_{j,m},y_{j,m}:j<\ell\}$.

\begin{hypothesis}[Uniform bootstrap determinant]\label{hyp:bootstrap-det}
For every $\ell\ge8$, the determinant
\[
  \Delta_\ell
  := \det A_\ell
\]
is not the zero complex-valued polynomial on these real lower-degree
coordinates.
\end{hypothesis}

\begin{lemma}[Bootstrap uniqueness]\label{lem:bootstrap}
Assume Hypothesis~\ref{hyp:bootstrap-det}. For every fixed
$\ell\ge8$, the set $\{\Delta_\ell\ne0\}$ is Zariski-open and
Lebesgue-co-null in the lower-degree coefficient space. On this set,
the entries in $\mathcal{T}_\ell$ uniquely determine
$\mathbf{F}_\ell$ from $(\mathbf{F}_0,\dots,\mathbf{F}_{\ell-1})$.
\end{lemma}
\begin{proof}
The entries of $A_\ell$ are polynomial functions of the known
real-coordinate lower degrees, and $b_\ell$ is read directly from the
selected complex bispectral entries. If $\Delta_\ell\ne0$, the square system
\eqref{eq:bootstrap-system} has the unique solution
$\mathbf{F}_\ell=A_\ell^{-1}b_\ell$ over $\C^{2\ell+1}$, hence also
on the real-signal slice. Since $\Delta_\ell$ is a nonzero
complex-valued polynomial by Hypothesis~\ref{hyp:bootstrap-det},
$\abs{\Delta_\ell}^2=(\re\Delta_\ell)^2+(\im\Delta_\ell)^2$ is a
nonzero real polynomial. Its zero set is therefore a proper algebraic
subvariety.
\end{proof}

\begin{remark}[What remains to prove uniformly]\label{rem:uniform-gap}
The determinant hypothesis is the exact Kakarala-style nonsingularity
condition for our selective recurrence. Existing floating-point runs
find $\Delta_\ell\ne0$ degree-by-degree for $\ell\le100$ at a
deterministic witness, but they do not by themselves prove a
finite-band theorem.
An unconditional all-$L$ proof requires a closed-form nonvanishing
identity for $\Delta_\ell$ (or an explicit all-$\ell$ witness). The
current document isolates that identity rather than replacing it by a
local Jacobian-rank claim.
\end{remark}

\subsection{Proof of Theorem~\ref{thm:main}}
\label{sec:proof}

\begin{proof}[Proof of Theorem~\ref{thm:main}]
Suppose $\Phiaug(f)=\Phiaug(f')$ for generic $f,f'$. By
Lemma~\ref{lem:gauge}, gauge-fix to $\tilde f,\tilde f'$ on the gauge
slice; $\SO(3)$-invariance preserves
$\Phiaug(\tilde f)=\Phiaug(\tilde f')$. It suffices to prove
$\tilde f=\tilde f'$.

\emph{Degrees $0$ and $1$.} By Lemma~\ref{lem:F0F1}, $A_0$ and $c$ are
determined, and $\mathbf{F}_0,\mathbf{F}_1$ agree.

\emph{Seed degrees.} By Hypothesis~\ref{hyp:seed}, the full seed block
uniquely determines
$(\mathbf{F}_2,\dots,\mathbf{F}_{L_{\mathrm{seed}}})$ on the gauge
slice, so $\tilde f,\tilde f'$ agree at all degrees
$\le L_{\mathrm{seed}}$.

\emph{Bootstrap degrees.} Inductively, suppose $\tilde f$ and
$\tilde f'$ agree at degrees $0,\dotsc,\ell{-}1$. Then the linear
system $A_\ell\mathbf{F}_\ell=b_\ell$ in
Section~\ref{sec:bootstrap} has identical coefficient matrix and
right-hand side for both signals. By Lemma~\ref{lem:bootstrap} and
(G5), $A_\ell$ is invertible. Hence
$\mathbf{F}_\ell=\mathbf{F}_\ell'$.

\emph{Conclusion.} Inducting up to $\ell=L$, $\tilde f=\tilde f'$, so
$f'\in\SO(3)\cdot f$. The set of generic signals is the
$\SO(3)$-saturation of the Zariski-open subset
$\mathcal{V}\subset\mathcal{G}$ cut out by (G1)--(G5)
(Definition~\ref{def:gen}), hence Lebesgue-co-null in $V_L$.
\end{proof}

\section{Output size matches the lower bound}
\label{sec:size}

\begin{proposition}[Lower bound]\label{prop:lower-bound}
Any continuous semialgebraic (in particular, polynomial) complete
$\SO(3)$-invariant map $V_L\to\R^k$ that separates the generic free
stratum must have
$k\ge (L+1)^2-\dim\SO(3)=(L+1)^2-3=\Omega(L^2)$.
\end{proposition}
\begin{proof}
The free generic stratum of the quotient $V_L/\SO(3)$ is a
semialgebraic manifold of dimension $(L+1)^2-3$. A continuous
semialgebraic injection into $\R^k$ cannot decrease semialgebraic
dimension, so $k$ must be at least this dimension.
\end{proof}

\begin{proposition}[Upper bound]\label{prop:size}
$|\Phiaug| = \Theta(L^2)$.
\end{proposition}
\begin{proof}
By construction (Definition~\ref{def:Phi-aug}), the full seed block has
at most $512$ complex bispectral scalars, independent of~$L$. For each
degree $\ell\ge8$, $\Sb$ allocates exactly $2\ell+1$ selected complex
bispectral entries, the rows of the square bootstrap system
\eqref{eq:bootstrap-system}. Thus
$\sum_{\ell=8}^L(2\ell+1)=O(L^2)$, so $|\Phiaug|=O(L^2)$. Any optional
CG-power diagnostics (Remark~\ref{rem:cg-optional}) preserve this
bound. The matching $\Omega(L^2)$ lower bound comes from
Proposition~\ref{prop:lower-bound}.
\end{proof}

\begin{table}[h]
  \caption{Asymptotic entry counts in $\Phiaug$. The full low-degree
    seed is constant-size; the high-degree selective part is linear
    per target degree.}
  \label{tab:entry-counts}
  \centering
  \small
  \begin{tabular}{cccc}
    \toprule
    degree range & selected $\beta$ & selected $P$ & total \\
    \midrule
    $\ell\le7$ & full seed, $\le512$ & optional & $O(1)$ \\
    $\ell\ge8$ & $2\ell+1$ & $O(\ell)$ & $O(\ell)$ \\
    \bottomrule
  \end{tabular}
\end{table}

\section{Computational evidence and limitations}
\label{sec:certs}

Theorem~\ref{thm:main} is structural once
Hypotheses~\ref{hyp:seed} and~\ref{hyp:bootstrap-det} are supplied.
The current repository does not yet contain either a finite-band proof
of Hypothesis~\ref{hyp:seed} or a closed-form proof of
Hypothesis~\ref{hyp:bootstrap-det} for all~$\ell$. It does contain
reproducible numerical evidence for the bootstrap determinants through
the band-limits used in the experiments.

\begin{description}[style=nextline]
  \item[\textnormal{C1.\ Finite seed injectivity (Hypothesis~\ref{hyp:seed}).}]
    Open in this document. Kakarala's theorem motivates the statement,
    but does not prove the finite-band scalar seed claim as written.
    A rigorous replacement would be a finite-band theorem or an
    exact/interval computation certifying the seed map.
  \item[\textnormal{C2.\ Bootstrap determinant evidence.}]
    Checked per-$\ell$ at the deterministic seed-42 witness for
    $\ell\in\{8,\dots,100\}$ in IEEE~754 double precision, with
    $\sigma_{\min}(A_\ell)/\sigma_{\max}(A_\ell) \ge 10^{-4}$
    in the existing runs. This is strong numerical evidence, not a
    mathematical certificate. The script is
    \texttt{verify\_linear\_bootstrap.py} in \texttt{benchmarks/}.
  \item[\textnormal{C3.}]
    \textbf{CG-power real-Jacobian diagnostics.} The augmented real Jacobian is numerically full-rank for
    $\ell\in\{4,\dots,30\}$ at the same witness. This is evidence that
    the implementation's augmented invariant is well-conditioned, but
    it is not used as the global uniqueness step in the corrected proof.
\end{description}

\begin{remark}[Finite band-limit evidence]\label{rem:numerical-evidence}
The floating-point checks through $\ell=100$ support the determinant
hypothesis for the band-limits used in the paper. They would become a
rigorous finite-band theorem only after being upgraded to exact
arithmetic, interval arithmetic, or a perturbation bound showing that
the exact matrix is certified away from singularity.
\end{remark}

\begin{remark}[The uniform determinant gap]
For $\ell\ge8$, the selected family $\mathcal{T}_\ell$ has a uniform
closed form, but we do not yet have a closed-form identity proving
$\Delta_\ell\not\equiv0$ for every $\ell$. Kakarala's full-bispectrum
recurrence avoids this issue by using full matrix Fourier blocks and
ordinary matrix inverses. Our selective scalar recurrence needs an
analogous nonvanishing identity for the determinant of the
$O(\ell)$-row matrix in~\eqref{eq:bootstrap-system}. Establishing that
identity is the remaining step needed to turn the conditional theorem
into an unconditional all-$L$ statement.
\end{remark}

\begin{remark}[Constructive seed alternative]
Earlier drafts used an explicit closed-form / finite-fibre /
degree-$4$ filter as the seed solver. That route is useful
algorithmically but depends on fibre enumeration and nonlinear
least-squares residual checks; it is deliberately not load-bearing in
the proof above. The theorem instead assumes finite seed injectivity
directly in Hypothesis~\ref{hyp:seed}.
\end{remark}

\section{Relationship to prior work}
\label{sec:related}

\begin{table}[h]
  \caption{Comparison of bispectral invariants for signals on $S^2$ or
    finite groups~$G$.}
  \label{tab:comparison}
  \centering
  \small
  \setlength{\tabcolsep}{4pt}
  \begin{tabular}{@{}lllc@{}}
    \toprule
    Result & Domain & Output size & Complete? \\
    \midrule
    \citet{kakarala2012bispectrum}
      & $S^2$, full bispec. & $O(L^3)$ & Yes \\
    \citet{edidin2024orbit}
      & $S^2$, full bispec.\ (band-lim.) & $O(L^3)$ & Yes (generic) \\
    \citet{bendory2025orbit}
      & $S^2$, frequency marching & $O(L^2)$/shell & Yes (3 shells) \\
    \citet{mataigne2024selective}
      & finite $G$, sel.\ bispec. & $O(\abs{G})$ & Yes (exact) \\
    \textbf{This work}
      & \boldmath$S^2$, sel.\ scalar bispec.\textsuperscript{$\dagger$}
      & $\boldsymbol{\Theta(L^2)}$ & Conditional; num.\ evidence \\
    \bottomrule
  \end{tabular}
  \par\smallskip
  {\footnotesize $\dagger$ Optional CG-power diagnostics may be
  appended (Remark~\ref{rem:cg-optional}); $\SP=\varnothing$ in the
  proof.}
\end{table}

\paragraph{Mataigne et al.\ (finite groups).}
\citet{mataigne2024selective} prove that for finite groups, an
$O(\abs{G})$-sized subset of the scalar bispectrum suffices for
completeness, by a breadth-first traversal of the irrep graph. Our
algorithm is the analog for $\SO(3)$ on $S^2$: selected bootstrap entries play the
role of edges, and the for-loop in
Algorithm~\ref{alg:so3-inversion} is a continuum BFS. Applied to a
finite group, our argument collapses to theirs --- the CG power
machinery becomes redundant because the finite-group reality
constraint is algebraic, and Lemma~\ref{lem:bootstrap} is replaced by
their closed-form Fourier-coefficient recurrence.

\paragraph{Kakarala (full bispectrum).}
\citet{kakarala1992triple,kakarala2012bispectrum} prove completeness
of the full $O(L^3)$-sized bispectrum on homogeneous spaces.
Our finite seed hypothesis is modeled on this full-bispectrum result,
but the theorem here does not derive finite-band seed injectivity from
Kakarala alone. The per-degree determinant bootstrap then amplifies
the assumed seed to higher degrees at $\Theta(L^2)$ output.

\paragraph{Edidin and Satriano (band-limited bispectrum).}
\citet{edidin2024orbit} prove generic separation by the full
band-limited bispectrum, with an ``up to reflection'' caveat
(Remark~\ref{rem:edidin}). The parity argument of
Corollary~\ref{cor:first-odd} explains why the apparent
orientation-reversing $T_R$ ambiguity is visible to the scalar
bispectrum once the seed includes degree~$4$.

\paragraph{Bendory et al.\ (frequency marching).}
\citet{bendory2025orbit} reconstruct band-limited spherical signals
from a frequency-marching scheme using $O(L^2)$ third-order moments
per shell. Our invariant is closer to the bispectral-coefficient
formulation of \citet{kakarala2012bispectrum} but achieves the same
asymptotic $\Theta(L^2)$ size.

\paragraph{Bispectral neural networks.}
\citet{sanborn2023bispectral} introduced bispectrum-based learnable
invariant pooling for finite groups; \citet{myers2025selective}
extend the selective-bispectrum construction to the disk. Our $\SO(3)$
conditional result is aimed at enabling the analogous architecture for
spherical signals.

\bibliographystyle{plainnat}
\bibliography{references}

\end{document}
